\documentclass[a4paper,11pt]{article}
\usepackage[utf8]{inputenc}
\usepackage[T1]{fontenc}
\usepackage[french]{babel}
\usepackage[left=2cm,right=2cm,top=2cm,bottom=2cm]{geometry}
\usepackage{xcolor}
\usepackage{hyperref}
\usepackage{fontawesome5}
\usepackage{parskip}

\definecolor{primary}{RGB}{0, 45, 100}
\hypersetup{colorlinks=true, urlcolor=primary}

\begin{document}

\noindent
\begin{minipage}[t]{0.5\textwidth}
    \textbf{Lo\"ic Aron MBASSI EWOLO}\\
    \'El\`eve-Ing\'enieur Bac+5 -- Double dipl\^ome ENSEA\\[3pt]
    \faMapMarker\ Cergy, France\\
    \faPhone\ (+33) 7 59 52 33 98\\
    \faEnvelope\ \href{mailto:loicaron.mbassiewolo@ensea.fr}{loicaron.mbassiewolo@ensea.fr}
\end{minipage}
\hfill
\begin{minipage}[t]{0.4\textwidth}
    \raggedleft
    Cergy, le 21 mars 2026
\end{minipage}

\vspace{1.2cm}

\hfill
\begin{minipage}[t]{0.5\textwidth}
    \raggedleft
    \textbf{Thales}\\
    Service Recrutement\\
    Rennes, France
\end{minipage}

\vspace{1cm}

\noindent\textbf{Objet :} Candidature en alternance -- Ing\'enieur Cybers\'ecurit\'e Logicielle \& IA -- Septembre 2026

\vspace{0.8cm}

Madame, Monsieur,

L'IA embarqu\'ee dans des syst\`emes critiques introduit des vecteurs d'attaque sp\'ecifiques que les m\'ethodologies de s\'ecurit\'e classiques ne couvrent pas : empoisonnement de mod\`eles, inversion, extraction par inf\'erence adversariale. Thales, leader en s\'ecurit\'e des syst\`emes de d\'efense et de s\'ecurit\'e, est l'un des rares acteurs \`a adresser ces menaces de fa\c{c}on syst\'ematique. Contribuer \`a l'\'evaluation de la s\'ecurit\'e logicielle de syst\`emes IA \`a Rennes s'inscrit directement dans mon double parcours IA/s\'ecurit\'e.

Mon stage \`a l'INRIA Saclay (\'equipe TRIBE, avril-ao\^ut 2026) m'a form\'e \`a l'analyse de risques NLP : identification de clauses contractuelles sensibles, scoring de conformit\'e RGPD, et \'etude de la robustesse de mod\`eles BERT face \`a des donn\'ees adversariales. En parall\`ele, mon projet Snappy (2025) d\'emontre une ma\^itrise concr\`ete de la s\'ecurit\'e applicative : impl\'ementation du protocole Signal (Double Ratchet), chiffrement AES-256-GCM, gestion des cl\'es de session dans une architecture microservices Spring WebFlux. Ces r\'ealisations couvrent les deux dimensions du poste -- IA et s\'ecurit\'e logicielle.

Mon projet PROJET-TAXI (C, IPC Linux, m\'emoire partag\'ee, signaux) illustre ma ma\^itrise des syst\`emes bas-niveau, compl\'ement\`ere \`a une expertise s\'ecurit\'e applicative. Ma formation ENSEA en syst\`emes embarqu\'es et ma formation en g\'enie logiciel \`a l'\'Ecole Polytechnique de Yaound\'e m'ont donn\'e une vision syst\`eme compl\`ete, utile pour l'audit de logiciels IA int\'egr\'es dans des environnements contraints.

Je suis disponible d\`es septembre 2026 pour cette alternance de 12 mois \`a Rennes. La perspective de contribuer \`a la s\'ecurit\'e de syst\`emes IA critiques avec des m\'ethodologies de niveau d\'efense constitue une opportunit\'e tr\`es motivante. Je reste \`a votre disposition pour un entretien.

Je vous prie d'agr\'eer, Madame, Monsieur, l'expression de mes salutations distingu\'ees.

\vspace{1cm}

\noindent\textbf{Lo\"ic Aron MBASSI EWOLO}

\end{document}
